% !TeX TS-program = xelatex
% 虚构演示简历 — 仅用于 hellojobs 测试数据
\documentclass{resume-photo}
\usepackage{xeCJK}
\setCJKmainfont{Noto Serif CJK SC}
\usepackage{fontspec}
\usepackage{enumitem}
\setlist[itemize]{itemsep=0pt, topsep=1pt, parsep=0pt, partopsep=0pt}

\ResumeName{唐莉}

\begin{document}

\ResumeContacts{
  138-0017-0017,%
  \ResumeUrl{mailto:tangli@tsinghua.edu.cn}{tangli@tsinghua.edu.cn},%
  \textnormal{清华大学 | 电子信息 · 硕士 | 2023—2026}%
}

\ResumeTitle

\noindent 推荐系统与排序模型，熟悉 PyTorch、DeepFM 与多任务学习；在离线 AUC 提升 0.018，线上 CTR +4.5\%。注重可观测性与文档沉淀，习惯在评审阶段拆解风险；参与线上复盘与跨团队联调，英语 CET-6，可阅读官方文档与 RFC（演示数据）。

\section{教育背景}
\ResumeItem
{清华大学}
[\textnormal{电子工程系 | 电子信息 | 硕士}]
[2023.09—2026.06(预计)]

\begin{itemize}
  \item 导师组：机器学习与推荐。
\end{itemize}


\ResumeItem
{学术交流与在线课程}
[\textnormal{暑校 / MOOC / 厂商认证（演示）}]
[2022—2025]

\begin{itemize}
  \item 完成《深入理解计算机系统》配套实验与 MIT 6.828 / 6.S081 公开课部分章节跟做。
  \item Coursera / edX 分布式系统、云计算相关专项课程结业证书 4 门（虚构编号）。
  \item 通过 AWS SAA / 阿里云 ACP / Redis 开发者等认证中的 1--2 项（简历测试数据）。
  \item 参加 CCF CCSP / 华为 ICT / 百度之星等学科竞赛集训营并完成全部作业。
\end{itemize}



\section{专业技能}
\begin{itemize}
  \item 熟悉 Python、Spark、特征工程与离线评估（AUC、GAUC）。
  \item 掌握 Wide\&Deep、DeepFM、DIN 等排序模型。
  \item 了解 TensorFlow Serving 与 A/B 实验平台对接。
  \item 熟悉 Linux 日常运维与 Shell，具备日志检索与 CPU/内存突增初步定位能力。
  \item 掌握 Git / CI 与 Code Review，了解 Docker 与 Kubernetes 基本编排与滚动发布。
  \item 了解 OAuth2 / JWT、接口幂等与 REST 规范，具备单元测试与集成测试习惯（JUnit/pytest 等）。
  \item 能阅读英文技术文档，具备跨团队沟通与需求澄清能力（测试简历）。
\end{itemize}

\section{实习经历}
\ResumeItem
{快手}
[\textnormal{推荐算法实习生}]
[2024.07—2024.12]

\begin{itemize}
  \item \textbf{职责}: 短视频粗排模型特征与负采样策略迭代。
  \item \textbf{成果}: 离线 Recall@100 提升 2.3\%，线上时长 +1.2\%。
\end{itemize}


\ResumeItem
{海康威视 | 行业应用}
[\textnormal{嵌入式 / 后端实习生}]
[2023.09—2024.02]

\begin{itemize}
  \item \textbf{设备接入}: 参与国标设备接入网关协议适配与异常码映射。
  \item \textbf{可靠性}: 实现断线重连与心跳退避策略，掉线恢复时间缩短 40\%。
  \item \textbf{日志}: 统一日志格式与 traceId 透传，问题定位时间减半。
  \item \textbf{客户}: 支持 2 次驻场联调，收集需求并转化为可验收用例。
  \item \textbf{规范}: 熟悉编码规范与静态检查门禁，CI 全绿方可合并。
\end{itemize}



\section{项目经历}
\ResumeItem
{多兴趣召回模块}
[\textnormal{科研助理项目}]
[2024.02—2024.08]

\begin{itemize}
  \item 胶囊网络思想做兴趣向量拆分，长尾作者曝光 +8\%。
\end{itemize}


\ResumeItem
{统一配置中心改造}
[\textnormal{Spring Cloud Config → Nacos（演示）}]
[2024.01—2024.04]

\begin{itemize}
  \item \textbf{目标}: 解决配置分散、无灰度与无审计问题，支撑 80+ 微服务。
  \item \textbf{方案}: 双写双读迁移，配置变更走审批流与版本快照。
  \item \textbf{结果}: 配置发布平均耗时从 25min 降至 3min，人为误配下降 60\%。
  \item \textbf{风险}: 设计回滚开关与只读兜底文件，演练 2 次成功。
\end{itemize}



\section{个人荣誉}
\begin{itemize}
  \item 清华大学 综合优秀奖学金
  \item Kaggle 竞赛 Notebook 银牌区
  \item 全国计算机等级考试四级（网络工程师）（演示）
  \item 校级「优秀学生干部 / 优秀团员」或技术社团骨干（综合测评前 15\%）
  \item 开源贡献：向知名项目提交被合并的 PR（文档/测试/feature）（演示）
\end{itemize}

\end{document}
